\Article{29}
Let \(t\) denote a given linear function of the quantities \(x,y,z\) etc.; say,
\[
t=fx+gy+hz+\text{etc.}+k
\]
The value of \(t\) obtained from the most plausible values of \(x,y,z\) etc. will therefore be \(=fA+gB+hC+\text{etc.}+k\), which we shall denote by \(K\). If this is adopted as the true value of \(t\), an error is committed, which will be
\SourceLeaf{040}
\[
=fEx+gEy+hEz+\text{etc.}
\]
and will be denoted by \(Et\). Clearly, the mean value of this error is \(=0\), or the error will be free of a constant part. But the mean value of the square \((Et)^2\), or the mean value of the sum
\[
\begin{array}{llll}
ff(Ex)^2&+2fg\,Ex.Ey&+2fh\,Ex.Ez&+\text{etc.}\\
&+gg(Ey)^2&+2gh\,Ey.Ez&+\text{etc.}\\
&&+hh(Ez)^2&+\text{etc. etc.}
\end{array}
\]
is, by what we set out in the preceding article, equal to the product of \(mmp\) and the sum
\[
\begin{array}{llll}
ff[\alpha\alpha]&+2fg[\alpha\beta]&+2fh[\alpha\gamma]&+\text{etc.}\\
&+gg[\beta\beta]&+2gh[\beta\gamma]&+\text{etc.}\\
&&+hh[\gamma\gamma]&+\text{etc. etc.}
\end{array}
\]
or to the product of \(mmp\) and the value of the function \(\Omega-M\) obtained by the substitutions
\[
\xi=f,\quad\eta=g,\quad\zeta=h\quad\text{etc.}
\]
Thus, denoting this particular value of the function \(\Omega-M\) by \(\omega\), the mean error to be feared while we adhere to the determination \(t=K\) will be \(=m\BareRoot{p\omega}\), or the weight of this determination will be \(=\frac1\omega\).

Since we have identically
\[
\Omega-M=(x-A)\xi+(y-B)\eta+(z-C)\zeta+\text{etc.}
\]
it is clear that \(\omega\) is also equal to the particular value of the expression
\[
(x-A)f+(y-B)g+(z-C)h+\text{etc.}
\]
or to the particular value of \(t-K\) obtained when the indeterminates \(x,y,z\) etc. are assigned those values that correspond to the values \(\xi,\eta,\zeta\) etc. being \(f,g,h\) etc.

Finally, we observe that if \(t\) is expressed identically as a function of \(\xi,\eta,\zeta\) etc., its constant part must be \(=K\). Thus, if identically
\[
t=F\xi+G\eta+H\zeta+\text{etc.}+K,\quad\text{then}\quad\omega=fF+gG+hH+\text{etc.}
\]
\SourceLeaf{041}
\Article{30}
The function \(\Omega\) attains its \textit{absolute minimum} value \(M\), as we saw above, on putting \(x=A,\allowbreak y=B,\allowbreak z=C\) etc., or \(\xi=0,\allowbreak\eta=0,\allowbreak\zeta=0\) etc. But if one of these quantities has already been assigned a \textit{different} value, e.g. \(x=A+\Delta\), then, as the others vary, \(\Omega\) can attain a \textit{relative minimum} value, which is clearly obtained by means of the equations
\[
x=A+\Delta,\quad\frac{d\Omega}{dy}=0,\quad\frac{d\Omega}{dz}=0\quad\text{etc.}
\]
We must therefore have \(\eta=0,\zeta=0\) etc., and hence the relation gives
\[
x=A+[\alpha\alpha]\xi+[\alpha\beta]\eta+[\alpha\gamma]\zeta+\text{etc.},\quad\xi=\frac{\Delta}{[\alpha\alpha]}
\]
At the same time we shall have
\[
y=B+\frac{[\alpha\beta]\Delta}{[\alpha\alpha]},\quad z=C+\frac{[\alpha\gamma]\Delta}{[\alpha\alpha]}\quad\text{etc.}
\]
The relative minimum value of \(\Omega\) is then \(=[\alpha\alpha]\xi\xi+M=M+\frac{\Delta\Delta}{[\alpha\alpha]}\). Conversely, we conclude from this that if the value of \(\Omega\) must not exceed the prescribed bound \(M+\mu\mu\), the value of \(x\) must necessarily lie between the bounds \(A-\mu\BareRoot{[\alpha\alpha]}\) and \(A+\mu\BareRoot{[\alpha\alpha]}\). It deserves notice that \(\mu\BareRoot{[\alpha\alpha]}\) is equal to the mean error to be feared in the most plausible value of \(x\) if we put \(\mu=m\BareRoot{p}\), that is, if \(\mu\) is equal to the mean error of observations to which the weight \(=1\) is assigned.

Let us investigate more generally the minimum value of \(\Omega\) that can occur for a given value of \(t\), where, as in the preceding article, \(t\) denotes the linear function \(fx+gy+hz+\text{etc.}+k\) and its most plausible value is \(=K\): let the prescribed value of \(t\) be denoted by \(K+\kappa\). From the theory of maxima and minima, it is known that the solution of the problem is to be sought from the equations
\[
\begin{aligned}
\frac{d\Omega}{dx}&=\theta\frac{dt}{dx}\\[3pt]
\frac{d\Omega}{dy}&=\theta\frac{dt}{dy}\\[3pt]
\frac{d\Omega}{dz}&=\theta\frac{dt}{dz}\quad\text{etc.}
\end{aligned}
\]
or \(\xi=\theta f,\allowbreak\eta=\theta g,\allowbreak\zeta=\theta h\) etc., where \(\theta\) denotes a multiplier that is still undetermined. Thus, if, as in the preceding article, we put \textit{identically}
\[
t=F\xi+G\eta+H\zeta+\text{etc.}+K
\]
we shall have
\SourceLeaf{042}
\[
\begin{aligned}
K+\kappa&=\theta(fF+gG+hH+\text{etc.})+K,\quad\text{or}\\[3pt]
\theta&=\frac{\kappa}{\omega}
\end{aligned}
\]
taking \(\omega\) in the same sense as in the preceding article. And since \(\Omega-M\) is identically a homogeneous function of the second degree in the indeterminates \(\xi,\eta,\zeta\) etc., it follows of itself that its value for \(\xi=\theta f,\allowbreak\eta=\theta g,\allowbreak\zeta=\theta h\) etc. becomes \(=\theta\theta\omega\), and hence that the minimum value that \(\Omega\) can attain for \(t=K+\kappa\) becomes \(=M+\theta\theta\omega=M+\frac{\kappa\kappa}{\omega}\). Conversely, if \(\Omega\) must not exceed a prescribed value \(M+\mu\mu\), the value of \(t\) must necessarily lie between the bounds \(K-\mu\BareRoot{\omega}\) and \(K+\mu\BareRoot{\omega}\), where \(\mu\BareRoot{\omega}\) is equal to the mean error to be feared in the most plausible determination of \(t\) if for \(\mu\) we take the mean error of observations to which the weight \(=1\) is assigned.

\clearpage
